% !TEX program = xelatex
\documentclass[12pt,a4paper]{article}

\usepackage[UTF8,heading=true,scheme=chinese]{ctex}
\renewcommand{\abstractname}{摘要}
\renewcommand{\contentsname}{目录}
\renewcommand{\refname}{参考文献}
\renewcommand{\figurename}{图}
\renewcommand{\tablename}{表}
\usepackage{geometry}
\geometry{a4paper,margin=2.5cm}
\usepackage{titling}
\usepackage{fancyhdr}
\usepackage{enumitem}
\usepackage{xcolor}
\usepackage{listings}
\usepackage{booktabs}
\usepackage{longtable}
\usepackage{tabularx}
\usepackage{array}
\usepackage{hyperref}
\hypersetup{
  colorlinks=true,
  linkcolor=black,
  citecolor=blue,
  urlcolor=blue,
  pdftitle={经管类本科建模与写作 AI 智能体辅助工作流说明书},
  pdfauthor={ALLBLUEUK}
}
\usepackage{graphicx}
\usepackage{caption}
\captionsetup{font=small,labelfont=bf}
\usepackage{tikz}
\usetikzlibrary{shapes,arrows.meta,positioning,fit,backgrounds}

\definecolor{shellbg}{HTML}{F5F5F5}
\definecolor{shellfg}{HTML}{1A1A1A}
\definecolor{prompt}{HTML}{0066CC}
\definecolor{comment}{HTML}{888888}
\definecolor{problembox}{HTML}{FFF5E6}
\definecolor{problemborder}{HTML}{D88E2C}
\definecolor{solutionbox}{HTML}{E8F4E8}
\definecolor{solutionborder}{HTML}{4A8C4A}

\lstdefinestyle{shell}{
  backgroundcolor=\color{shellbg},
  basicstyle=\ttfamily\small\color{shellfg},
  frame=single,
  framerule=0.4pt,
  rulecolor=\color{black!30},
  breaklines=true,
  breakatwhitespace=true,
  columns=fullflexible,
  showstringspaces=false,
  keepspaces=true,
  literate={→}{{$\rightarrow$}}1
           {⏸}{{$\blacksquare$}}1
           {✓}{{\checkmark}}1
           {⚠}{{!}}1
           {…}{{\ldots}}1,
}
\lstdefinestyle{prompt}{
  style=shell,
  basicstyle=\ttfamily\small\color{prompt}\bfseries,
  backgroundcolor=\color{shellbg},
}
\lstdefinestyle{md}{
  backgroundcolor=\color{shellbg},
  basicstyle=\ttfamily\footnotesize\color{shellfg},
  frame=single,
  framerule=0.4pt,
  breaklines=true,
  columns=fullflexible,
  showstringspaces=false,
  keepspaces=true,
}

% Problem & solution boxes
\usepackage[most]{tcolorbox}
\newtcolorbox{problembox}[1][]{
  colback=problembox,colframe=problemborder,
  fonttitle=\bfseries,
  title=#1
}
\newtcolorbox{solutionbox}[1][]{
  colback=solutionbox,colframe=solutionborder,
  fonttitle=\bfseries,
  title=#1
}

\title{\Huge 经管类本科建模与写作\\
AI 智能体辅助工作流说明书\\
\large 实际操作使用手册 \quad{} ueaiworkflow}
\author{ALLBLUEUK\\
\url{https://github.com/ALLBLUEUK/ueaiworkflow}}
\date{2026 年 5 月}

\setlength{\headheight}{13.6pt}
\pagestyle{fancy}
\fancyhf{}
\fancyhead[L]{\small ueaiworkflow — 操作使用手册}
\fancyhead[R]{\small \thepage}

\begin{document}

\maketitle
\thispagestyle{empty}

\begin{abstract}
\noindent
本工作流不是孤立的技术产品，而是两年实际教学中针对经管类本科学生反复出现的具体困难沉淀的一套\textbf{可被 AI 智能体执行}的教学支架。

本说明书围绕一个核心问题展开：\emph{在数字经济与 AI 背景下，本科学生面对一个真实经贸问题时，如何从最初的兴趣点一路推进到可提交、可核验、可归档的研究报告，并在这个过程中既用上 AI 工具又不被 AI 工具替代？}

\textbf{本工作流的定位是``辅助学习'', 不是``替代研究''。} 学生自己的工作（提出问题、跑数据库、本机跑 GEMPACK / Stata、独立判断、写论文正文）始终由学生完成；AI 提供的是\textbf{学习理解的脚手架}: 把 UN Comtrade 字段翻译成中文、把 \texttt{.cmf} 每一行解释清楚、把 ViewSOL 里的数字翻译成方向 / 量级 / 经济含义三件套、按五维度检查论文结构是否塌。

读者读完本说明书将清楚: 项目在课堂中遇到了哪些具体问题、本工作流为这些问题提供了哪些具体回应、教师与学生如何在自己的电脑上使用它、以及一次端到端运行的完整输入与输出。
\end{abstract}

\tableofcontents
\newpage

%--------------------------------------------------------------------
\section{本工作流要解决的项目问题}
%--------------------------------------------------------------------

本工作流以《贸易数据库与分析工具》和《经济建模与写作》两门课程为支点，目标是把数据理解、模型操作、研究表达与 AI 工具组织成一条完整的能力链条。教学过程中反复观察到以下五类具体困难：

\begin{problembox}[问题 1：学生面对真实经贸问题时同时被五件事压住]
本科生第一次面对``关税冲击对中国电子部门的影响''这种真实问题时，要同时完成：
\begin{enumerate}[leftmargin=*]
  \item 从宽泛兴趣中收窄出一个可研究的问题；
  \item 在 UN Comtrade / WITS / GTAP base 中检索数据，并核对 HS 编码、单位、年份；
  \item 跑 GTAP/GEMPACK 或写回归代码，处理报错；
  \item 把模型输出读成方向 + 量级 + 经济含义；
  \item 把以上全部组织成一篇带文献综述、结论与边界的论文。
\end{enumerate}
\textbf{后果}：学生通常在第 1–2 步停滞或越级直接跳到第 5 步写文本，缺中间证据；少数勉强跑通的学生也无法说明每一步的来源。
\end{problembox}

\begin{problembox}[问题 2：AI 工具不规范使用容易退化为``替写'']
如果不设结构，学生使用 AI 极易停在以下三种浅层模式：
\begin{itemize}[leftmargin=*]
  \item 直接让 AI 写整段论文，自己只做润色；
  \item 让 AI 列出 HS 编码、GTAP set、回归系数等\emph{专业知识}，却不回到数据库 / 软件 / 手册核验；
  \item 把数据 / 模型 / 文献 / 写作四个环节混在一个会话里，最后无法区分哪些是 AI 的、哪些是自己的。
\end{itemize}
\textbf{后果}：AI 看似提高了效率，实际上削弱了学生的独立判断；教师评阅时也无法判断学生掌握了什么。
\end{problembox}

\begin{problembox}[问题 3：专业工具门槛高，挫败率集中在最初几小时]
GTAP 11、GEMPACK 命令、\texttt{.cmf} 文件、\texttt{.har} 头、Armington closure 这些专业概念对本科生属于``一旦走错就找不到入口''的陡坡。一份 \texttt{tariffshock.cmf} 写错路径或 \texttt{Steps} 设小了，会让学生在第一周耗去三四个小时——而问题往往只是工作目录或细分步数。
\end{problembox}

\begin{problembox}[问题 4：教师评阅缺乏过程证据，质量难以保障]
当学生只提交一份 PDF 时，教师无法判断：数据是否真的从 UN Comtrade 导出？模型是否本机跑通？回归是否做了平行趋势检验？AI 修改了多少？这些过程性证据如果不在课程结构里强制收集，期末就只剩文本结果可看，评阅的判断力下降。
\end{problembox}

\begin{problembox}[问题 5：项目结项时，过程材料散落各处]
项目要把课程改革沉淀为可结项、可推广的成果，但学生作业、小组报告、AI 使用记录、模型文件、最终论文如果分散在邮箱、QQ 群、各自电脑上，结项时无法快速汇总成佐证材料。
\end{problembox}

\medskip
本工作流回应这五个问题。下面三节依次说明：\textbf{要干什么}（第 \ref{sec:goal} 节）、\textbf{怎么干}（第 \ref{sec:six-stage} 节）、\textbf{在电脑上具体怎么操作}（第 \ref{sec:install} 与 \ref{sec:walkthrough} 节）。

%--------------------------------------------------------------------
\section{工作流要做的事}\label{sec:goal}
%--------------------------------------------------------------------

本工作流的核心目的是：\textbf{把一个真实经贸问题转化为六个分阶段、互相衔接、各有 AI 帮手、各有学生人工核验段、各有归档物证的课程任务}，让学生在 AI 协助下走完一次完整的研究过程，而教师在期末能拿到可评、可证的过程材料。

具体要完成的任务可以拆为五件：

\begin{solutionbox}[任务 1：把一个真实问题切成 6 个可执行阶段]
工作流强制把每个课程项目拆成 \texttt{00\_problem} $\rightarrow$ \texttt{01\_data} $\rightarrow$ \texttt{02\_model} $\rightarrow$ \texttt{03\_results} $\rightarrow$ \texttt{04\_writing} $\rightarrow$ \texttt{05\_feedback} 六阶段，每阶段对应一个独立 markdown 文件，学生与 AI 共同填写。学生在任何一个时刻打开项目，都能立刻看见自己走到哪里、下一步该做什么。
\end{solutionbox}

\begin{solutionbox}[任务 2：在每个阶段给 AI 设具体的脚本与边界]
每阶段配套一个 \emph{技能}（位于 \texttt{.claude/skills/}）和若干\emph{提示词模板}（位于 \texttt{prompts/}）。AI 不会面对一个空白上下文乱写——它读到``这是 stage 01，调用 \texttt{ueai-data-path} 技能，依据 \texttt{prompts/data-path-comtrade.md} 给出 UN Comtrade 字段名、HS 编码、检索 URL、核验清单''——AI 的产出因此结构化、可重复、可比较。
\end{solutionbox}

\begin{solutionbox}[任务 3：强制学生在每个阶段做人工核验]
每个阶段文件都有两段必填区：\texttt{\#\# AI 辅助记录}（由 AI 写）与 \texttt{\#\# 人工核验与修改}（学生必填）。AI 在每阶段结束\emph{暂停}并明确提示学生回到 UN Comtrade 截图核对、跑一遍 \texttt{runGTAP.bat}、对照教师讲义。学生不补齐人工核验段，\texttt{/ueai-validate} 命令会报失败。
\end{solutionbox}

\begin{solutionbox}[任务 4：把数据 / 模型 / AI 使用过程沉淀为可归档证据]
每个项目自带 \texttt{evidence/} 子目录，按 \texttt{data/} \texttt{model/} \texttt{results/} \texttt{paper/} 分类存放学生导出的原始 CSV、GTAP \texttt{.sl4}/\texttt{.har}、ViewSOL 截图、最终 PDF。一份自动生成的 \texttt{evidence\_checklist.md} 列出 7 项必备证据，\texttt{/ueai-archive} 命令逐项核对状态。
\end{solutionbox}

\begin{solutionbox}[任务 5：把课程作业、报告、论文与项目结项绑定]
工作流自带 \texttt{rules/report-quality-rubric.md}（五维度评分量规）与 \texttt{rules/academic-integrity.md}（学术诚信与 AI 披露规则）。教师只要给学生分发仓库 URL，全班作业、报告、期末论文的目录结构、AI 使用披露、评阅记录就统一了。学期末教师把每位学生的 \texttt{<slug>/} 文件夹打包，即可作为课程档案、结项佐证、推广素材。
\end{solutionbox}

\medskip
换句话说，本工作流面向的不是``写一个 Python 包''或``给学生一个 ChatGPT 网址''这种点状任务，而是把项目主题 ``\emph{数字经济与 AI 背景下经管类课程能力链条构建}'' 转化为可在每一台学生电脑上跑通的、由 AI 智能体执行的具体教学动作。

%--------------------------------------------------------------------
\section{六阶段：每个阶段在解决一个具体困难}\label{sec:six-stage}
%--------------------------------------------------------------------

下面逐阶段说明：\textbf{这个阶段在解决学生遇到的什么困难}、\textbf{它要求 AI 做什么}、\textbf{它要求学生做什么}、\textbf{它产出什么文件}。

\subsection{阶段 00 — 问题拆解（\texttt{00\_problem/problem\_brief.md}）}

\paragraph{解决什么困难。}学生说``我想做关税战''——一个泛主题、无国家、无时间、无变量、无方法。如果不强行拆解，学生大概率写出一篇``概述性''论文，无法继续推进。

\paragraph{AI 做什么。}AI 调用 \texttt{ueai-problem-framing} 技能，依据 \texttt{prompts/problem-framing.md} 把题目拆成五个字段——研究对象 / 核心变量 / 可用数据 / 可能方法 / 结果展示——并追加 3–5 条``学生必须独立核验的口径''。AI 不替学生选定方法，只给两条候选并标注假设。

\paragraph{学生做什么。}学生回到课堂讲义与数据库门户独立核验那 3–5 条口径，并在 \texttt{\#\# 人工核验与修改} 段写下教师要求、数据可得性、本人取舍。

\paragraph{产出。}一份明确陈述``中国对美电子部门 / 2017–2024 / 25pp 关税冲击 / GTAP 主路径 + DID 备选''的 \texttt{problem\_brief.md}，可作为之后所有阶段的研究边界。

\subsection{阶段 01 — 数据路径（\texttt{01\_data/data\_plan.md}）}

\paragraph{解决什么困难。}学生面对 UN Comtrade 网站不知道勾哪些字段；HS 编码三位还是六位选不准；USTR Section 301 List 1–4 的生效日与覆盖范围记不清；GTAP database 11 的 ``Electronics'' 部门到底包不包 HS 84。

\paragraph{AI 做什么。}AI 调用 \texttt{ueai-data-path} 技能，依据 \texttt{prompts/data-path-comtrade.md} 与 \texttt{prompts/data-path-io-table.md} 给出：主数据源（Reporter / Partner / Trade Flow / Classification / Cmd Code / Frequency / Period 七字段全填）、辅助数据源（WITS 关税、USTR 清单）、模型基年（GTAP 11 聚合方案）、一个可点击的 UN Comtrade 检索 URL 模板，以及 7 项必须线下核验的口径。

\paragraph{学生做什么。}学生在 UN Comtrade Plus 中按字段实际检索、导出 \texttt{.csv}、截图，放进 \texttt{evidence/data/}；并在 \texttt{\#\# 人工核验与修改} 写明实际选用的 HS 编码与基年汇率。

\paragraph{产出。}一份可执行的数据检索方案 + 一组真实导出的数据文件，确保后续模型阶段不会因为数据口径问题反工。

\subsection{阶段 02 — 模型与工具（\texttt{02\_model/model\_notes.md}）}

\paragraph{解决什么困难。}本阶段是 GTAP/GEMPACK 课程典型的``入门陡坡''。学生第一次看见 \texttt{tariffshock.cmf} 文件不知道每一行什么意思；第一次看到 \texttt{*** Newton convergence failed} 不知道是数据问题还是冲击太大；第一次看到 ViewSOL 里 100+ 个变量不知道哪些该看。任何一个细节卡住都可能让当周作业掉队。

\paragraph{AI 做什么。}AI 调用 \texttt{ueai-model-explain} 技能，输出\textbf{逐行解释式的学习材料}，包括：

\begin{enumerate}[leftmargin=*]
  \item 带中文逐行注释的 \texttt{.cmf} 模板。例如：
\begin{lstlisting}[style=md,basicstyle=\ttfamily\scriptsize]
File (new) tariff25pp.sl4 ;
! 求解输出文件；扩展名 .sl4 = "solution level 4"，
! 含百分比变化结果，用 ViewSOL 打开

File rwsbase = D:\GEMPACK12\GTAP\v11\base12x4.har ;
! base 数据文件；rwsbase 是 GTAP 标准模型固定的关键字
! .har 文件是 GEMPACK 的 header-array 二进制格式

Method = Gragg ;
Steps = 3 6 9 ;
Subintervals = 3 ;
! Gragg 多步法求解；Steps 3 6 9 意为分三次细分迭代
! 数字越大越精确，求解越慢。25pp 大冲击建议 6 9 12

Shock tms("ELE", "CHN", "USA") = 25 ;
! tms 是 import tariff power；语法 tms(部门,源国,目国)
! = 25 含义为提升 25 个百分点（不是 25%）
\end{lstlisting}

  \item \textbf{报错诊断手册}。每条报错给出原因 + 修复 + 怎么防止：
\begin{lstlisting}[style=md,basicstyle=\ttfamily\scriptsize]
报错：*** Could not open file rwsbase.har
原因：.cmf 中相对路径解析失败，工作目录不在 base 文件夹
修复：把 File rwsbase = base12x4.har ;
      改为 File rwsbase = D:\GEMPACK12\GTAP\v11\base12x4.har ;
防止：始终在 .cmf 中用绝对路径
\end{lstlisting}

  \item \textbf{ViewSOL 阅读地图}：告诉学生 GTAP 100+ 输出变量中应该重点看 \texttt{qgdp}（实际 GDP）、\texttt{EV}（等价变差）、\texttt{qxw}（出口数量）、\texttt{pxw}（出口价格）、\texttt{tot}（贸易条件），每个变量的单位（\% 还是百万美元）、按什么集合（区域 / 部门 / 双边）展开。

  \item \textbf{模型五条核心假设的明示}：完全竞争 + Armington 替代弹性 + 关税收入归代表性家庭 + 标准短期 closure + 投资按 RoRC 分配——并用一句话说明每条假设对结果的影响方向。

  \item 若学生暂无 GEMPACK 访问，AI 同时给出 Stata DID 备选代码骨架（变量定义、固定效应、聚类标准误、平行趋势事件研究图）。
\end{enumerate}

\paragraph{学生做什么。}学生在本机 PowerShell 中实际跑 \texttt{runGTAP.bat}（或 \texttt{stata did.do}），把 \texttt{.sl4}、\texttt{.har}、\texttt{.log} 放入 \texttt{evidence/model/}，并在 \texttt{\#\# 人工核验与修改} 写明本地路径、closure 选择、教师补充。

\paragraph{产出。}一份\textbf{可被学生读懂}的模型运行笔记。学生在跑 GEMPACK 之前看完这份笔记，应能回答：``为什么 Steps 设为 3 6 9''、``EV 单位是什么''、``如果 Newton 收敛失败我该怎么调参''。

\subsection{阶段 03 — 结果解释（\texttt{03\_results/result\_interpretation.md}）}

\paragraph{解决什么困难。}学生看见 ViewSOL 里的 \texttt{qgdp = -0.34\%}，常常只会复述``GDP 下降''——既不报量级，也不解释经济含义，更不与文献比较；或者反过来用``显著影响''``重大变化''这类无量词形容词代替数字。

\paragraph{AI 做什么。}AI 调用 \texttt{ueai-result-interpret} 技能，依据 \texttt{prompts/result-interpret.md} 把每个数值翻译为``\textbf{方向 + 量级 + 经济含义}''三件套（``-0.34\% 符号符合理论；以 2017 中国 GDP 为分母对应约 418 亿美元；损失主要由电子部门下游溢出''）。同时给出受损前 3 / 受益前 3 部门、4 类稳健性边界、与 Caliendo–Parro / Amiti et al. / Itakura–Lee 的量级对照。\textbf{所有数字强制带 \texttt{[来源: evidence/results/...]} 标记，} AI 若不确定数字则写 TODO 而不伪造。

\paragraph{学生做什么。}学生用本地求解的真实数值替换 AI 列出的占位，确认每条 \texttt{[来源: ...]} 标记指向真实文件；并在 \texttt{\#\# 人工核验与修改} 说明哪些数值与 AI 量级带不一致、原因为何。

\paragraph{产出。}一份``每条数值都带来源''的结果解释，可被教师任意抽查复核。

\subsection{阶段 04 — 论文与报告（\texttt{04\_writing/paper\_outline.md}）}

\paragraph{解决什么困难。}本科生写课程论文最容易塌的不是文笔，而是结构：研究问题模糊、文献堆砌无 gap、数据方法两段合并写不清识别假设、结果段堆数字无解释、结论段过度推论无边界。

\paragraph{AI 做什么。}AI 调用 \texttt{ueai-paper-structure} 技能，依据 \texttt{prompts/paper-structure.md} 与 \texttt{rules/report-quality-rubric.md} 给出五维度大纲：研究问题 / 文献综述 / 数据方法 / 结果 / 结论与边界。每个 §4 数值都引回 \texttt{03\_results/...} 并保留 \texttt{[来源: ...]} 标记；§5 结论必须含一段``本研究不能说明……''；论文末尾必须含 \texttt{\#\# 学术诚信声明} 段。\textbf{AI 不写任何完整正文段落，只给大纲与逐条要点。}

\paragraph{学生做什么。}学生在大纲基础上独立撰写正文（摘要、引言、文献综述、结论），补足中文参考文献 $\geq$ 4 条，并填学术诚信声明中的``本人完成的具体环节''。

\paragraph{产出。}一份结构清晰、有据可查、含 AI 使用披露的论文骨架，配套学生本人撰写的正文，最终生成 \texttt{evidence/paper/final.pdf}。

\subsection{阶段 05 — AI 反馈整理（\texttt{05\_feedback/ai\_revision\_log.md}）}

\paragraph{解决什么困难。}学生学期末交一份漂亮 PDF，教师无法判断 AI 在每一步的实际贡献——是 90\% 替写还是 10\% 检查？

\paragraph{AI 做什么。}AI 扫描 stage 00–04 全部 \texttt{\#\# AI 辅助记录} 段，整理成一张总览表（阶段 / 模型 / 提示词来源 / 接受–修改–拒绝状态）+ 五段详细记录。本阶段 AI 不生成新建议，只做日志整理。

\paragraph{学生做什么。}学生校对状态标注、补足``修改字数''、确认所有 \texttt{[来源: ...]} 与学术诚信声明在位。

\paragraph{产出。}一份审计级别的 AI 使用日志，让教师 / 评阅 / 项目结项审查可以快速判断学生与 AI 各自的贡献。

\subsection{评阅与归档命令}

除六阶段外，三个工具命令把工作流闭环：

\begin{itemize}[leftmargin=*]
  \item \texttt{/ueai-review} —— 切换到评阅人格（数据评阅 / 模型评阅 / 论文评阅），按 \texttt{rules/report-quality-rubric.md} 五维度打分，把结果\textbf{追加}到 \texttt{05\_feedback/ai\_revision\_log.md} 末尾，不修改论文本身。
  \item \texttt{/ueai-archive} —— 遍历 \texttt{evidence\_checklist.md} 7 项，依据 \texttt{evidence/} 实际文件存在情况自动勾选 \texttt{[x]} / \texttt{[\textasciitilde]} / \texttt{[\,]}，列出缺失项与建议补充。
  \item \texttt{/ueai-validate} —— 提交前最终检查：六阶段文件齐全 / AI 段非空 / 人工核验段非空 / 证据 $\geq$ 5/7 / 所有数字带来源 / 含学术诚信声明 / AI 反馈日志存在。八项全过给 \texttt{READY FOR SUBMISSION}；否则列出待修补项。
\end{itemize}

\medskip
六阶段 + 三工具命令的完整流水线如下图：

\begin{center}
\begin{tikzpicture}[
  node distance=0.7cm and 0.5cm,
  stage/.style={draw,rounded corners,inner sep=4pt,minimum width=2.0cm,minimum height=0.95cm,align=center,font=\footnotesize},
  tool/.style={draw,rounded corners,inner sep=4pt,minimum width=2.0cm,minimum height=0.95cm,align=center,font=\footnotesize,dashed},
  arrow/.style={-Latex,thick,gray}
]
\node[stage,fill=blue!15]  (s0) {00\\\textbf{问题}};
\node[stage,fill=blue!15,right=of s0]  (s1) {01\\\textbf{数据}};
\node[stage,fill=blue!15,right=of s1]  (s2) {02\\\textbf{模型}};
\node[stage,fill=blue!15,right=of s2]  (s3) {03\\\textbf{结果}};
\node[stage,fill=blue!15,right=of s3]  (s4) {04\\\textbf{写作}};
\node[stage,fill=blue!15,right=of s4]  (s5) {05\\\textbf{反馈}};
\node[tool,fill=yellow!20,below=1.4cm of s2,xshift=-0.5cm] (rv) {/ueai-review};
\node[tool,fill=yellow!20,below=1.4cm of s3,xshift=0.0cm] (ar) {/ueai-archive};
\node[tool,fill=yellow!20,below=1.4cm of s4,xshift=0.5cm] (va) {/ueai-validate};
\draw[arrow] (s0) -- (s1);
\draw[arrow] (s1) -- (s2);
\draw[arrow] (s2) -- (s3);
\draw[arrow] (s3) -- (s4);
\draw[arrow] (s4) -- (s5);
\draw[arrow,dashed] (s5) -- (rv);
\draw[arrow,dashed] (s5) -- (ar);
\draw[arrow,dashed] (s5) -- (va);
\end{tikzpicture}
\end{center}

%--------------------------------------------------------------------
\section{安装与启动}\label{sec:install}
%--------------------------------------------------------------------

工作流不需要 Python、不需要 \texttt{pip install}、不需要构建。仓库本身就是配置——克隆下来 AI 智能体直接读取。

\subsection{适用的 AI 客户端}

教师与学生可在以下任何一种 AI 客户端中使用本工作流；所有客户端共享同一份 \texttt{AGENTS.md} 作为指令入口：

\begin{itemize}[leftmargin=*]
  \item \textbf{Claude Code}（推荐）—— 启动后自动读取 \texttt{AGENTS.md} 与 \texttt{.claude/}，直接输入 \texttt{/ueai-new} 启动。
  \item \textbf{Codex CLI} —— 原生支持 \texttt{AGENTS.md}，自然语言触发。
  \item \textbf{Cursor} —— 把 \texttt{AGENTS.md} 加为 ``Rules''。
  \item \textbf{GitHub Copilot Chat}（VS Code） —— 仓库自带 \texttt{.github/copilot-instructions.md}。
  \item \textbf{Aider / Continue} —— 把 \texttt{AGENTS.md} 加为 read-only context。
  \item \textbf{无 AI 客户端} —— 也可手动复制 \texttt{templates/} 与 \texttt{prompts/}，粘到任意 LLM 网页版调用。
\end{itemize}

\subsection{操作步骤}

\begin{lstlisting}[style=shell,caption={第一次使用（任意学生电脑，PowerShell）}]
PS> cd D:\courses\trade-db\2026-spring\student-007
PS> git clone https://github.com/ALLBLUEUK/cuibe-uea-workflow .ueai
PS> claude        # 或 codex / cursor / vscode
\end{lstlisting}

启动 AI 助手后，按下述顺序输入命令：

\begin{lstlisting}[style=prompt]
> /ueai-new
> 课程：贸易数据库与分析工具
> 题目：美国加征关税的宏观影响分析
> 研究问题：2018 年以来美国对华加征关税对中国电子部门出口与宏观福利的影响如何
> slug: tariff-impact

> /ueai-stage 00     # 问题拆解
> /ueai-stage 01     # 数据路径
> /ueai-stage 02     # 模型与工具
> /ueai-stage 03     # 结果解释
> /ueai-stage 04     # 论文与报告
> /ueai-stage 05     # AI 反馈整理

> /ueai-review       # 教师 / 同伴评阅
> /ueai-archive      # 归档清单
> /ueai-validate     # 提交前检查
\end{lstlisting}

每个 \texttt{/ueai-stage} 命令结束时 AI 会暂停，明确告诉学生``请在 PowerShell 中跑 X''或``请在 UN Comtrade 截图 Y''——学生完成后回到会话再输入下一个命令。完整 PowerShell 输入与 AI 输出见第 \ref{sec:walkthrough} 节。

%--------------------------------------------------------------------
\section{端到端实际操作记录}\label{sec:walkthrough}
%--------------------------------------------------------------------

\textbf{运行环境}：Windows 10 Enterprise 22H2 + PowerShell 7.4 + Windows Terminal 1.18 + Claude Code 2.0 (Claude Opus 4.7)。

\textbf{示范项目}：\texttt{tariff-impact} — 美国 25 个百分点关税对中国电子部门冲击的 GTAP 模拟（项目仓库 \texttt{examples/tariff-impact/} 完整保存了下述每一步的输出文件）。

\textbf{读法}：每节标题对应一次用户输入；用户输入用 \emph{蓝色加粗等宽}显示；AI 输出节录与实际写入文件的内容随后给出。

\subsection{第 0 步：克隆 + 启动}

\begin{lstlisting}[style=shell]
PS D:\courses\trade-db\2026-spring\student-007>
  git clone https://github.com/ALLBLUEUK/ueaiworkflow .ueai

Cloning into '.ueai'...
remote: Counting objects: 78, done.
Receiving objects: 100% (78/78), 142.31 KiB | 1.20 MiB/s, done.

PS D:\courses\trade-db\2026-spring\student-007>
  claude        # 启动 Claude Code，自动读取 .ueai/AGENTS.md
\end{lstlisting}

\subsection{第 1 步：/ueai-new}

\noindent\textbf{用户输入}：
\begin{lstlisting}[style=prompt]
> /ueai-new
> 课程：贸易数据库与分析工具
> 题目：美国加征关税的宏观影响分析（电子部门）
> 研究问题：2018 年以来美国对华加征关税对中国电子部门出口与宏观福利的影响如何
> slug: tariff-impact
\end{lstlisting}

\noindent\textbf{AI 输出（实际写入磁盘 = 8 个 .md + 4 个 evidence/ 子目录）}：
\begin{lstlisting}[style=shell]
读取了 AGENTS.md §4 与 .claude/commands/ueai-new.md。

✓ Created tariff-impact/
  |-- README.md
  |-- 00_problem/problem_brief.md
  |-- 01_data/data_plan.md
  |-- 02_model/model_notes.md
  |-- 03_results/result_interpretation.md
  |-- 04_writing/paper_outline.md
  |-- 05_feedback/ai_revision_log.md
  |-- evidence/
  |   |-- evidence_checklist.md
  |   |-- data/      (空)
  |   |-- model/     (空)
  |   |-- results/   (空)
  |   `-- paper/     (空)
  `-- prompt_templates.md

所有占位符已替换。AI 内容尚未填入 - 运行 /ueai-stage 00 开始第一阶段。
\end{lstlisting}

\subsection{第 2 步：/ueai-stage 00 — 问题拆解}

\noindent\textbf{用户输入}：
\begin{lstlisting}[style=prompt]
> /ueai-stage 00
\end{lstlisting}

\noindent\textbf{AI 实际写入 \texttt{00\_problem/problem\_brief.md} 的核心内容（节选）}：
\begin{lstlisting}[style=md]
研究对象：中国对美出口（货物贸易），HS 2017 第 84、85 章；
时间窗 2017—2024 年；以 2017 为基年（GTAP 11 基年）。

核心变量：
- 被解释变量：中国对美电子产品出口额（USD nominal）、
  中国 GDP（实际，量价分解后的 qgdp）、
  中国福利（等价变差 EV，2017 基价亿美元）。
- 解释变量 / 冲击：美国对华关税增量 Delta-tau（百分点）、
  HS 子章节加征覆盖率（0–1）、上游中间品依存度。

可用数据：
- UN Comtrade Plus（双边贸易，HS 2017 6 位）
- WITS Tariff（MFN + AVE for Section 301）
- USTR Section 301 List 1–4 的 HS 编码清单
- GTAP database 11（结构性基准，2017 base）

可能方法：
(A) GTAP / GEMPACK 反事实模拟：设 25pp 关税冲击于 Electronics 部门
    （中→美），求解 EV、qgdp、qxw；
    假设完全竞争 + Armington + 关税收入归代表性家庭。
(B) 行业 DID（备选）：以 USTR 清单覆盖的 HS6 为处理组，
    非覆盖 HS6 为控制组，2018 年为事件年；
    面板模型，固定效应 + 聚类标准误。

可核验口径（学生必须独立完成）：
- [ ] HS 8542 (集成电路)、8517 (通信)、8471 (计算机) 是否覆盖
- [ ] USTR Section 301 List 1–4 各自生效日与 HS 映射
- [ ] UN Comtrade 镜像核验
- [ ] 2024 年度数据可得性
- [ ] GTAP 11 的 Electronics 部门口径
\end{lstlisting}

\noindent\textbf{AI 暂停提示}：
\begin{lstlisting}[style=shell]
[Pause] 请在 00_problem/problem_brief.md 的「## 人工核验与修改」段填写：
   1) 数据来源核验（UN Comtrade 检索截图）
   2) 教师要求对照（课堂讲义中的部门聚合口径）
   3) 修改说明
填好后运行 /ueai-stage 01。
\end{lstlisting}

\subsection{第 3 步：/ueai-stage 01 — 数据路径}

\noindent\textbf{用户输入}：
\begin{lstlisting}[style=prompt]
> /ueai-stage 01
\end{lstlisting}

\noindent\textbf{AI 实际写入 \texttt{01\_data/data\_plan.md} 的字段表（节选）}：
\begin{lstlisting}[style=md]
A. 主数据源：UN Comtrade Plus（年度、HS 2017）

| 字段           | 取值                                              |
|----------------|---------------------------------------------------|
| Reporter       | 156 (China)                                       |
| Partner        | 842 (USA)                                         |
| Trade Flow     | X (Export) 与 M (Import) 镜像核验                 |
| Classification | HS 2017                                           |
| Cmd Code       | 8542, 8517, 8471, 8473, 8504, 8528                |
| Frequency      | A (年度)                                          |
| Period         | 2017–2024                                         |

检索 URL 示例：
https://comtradeplus.un.org/TradeFlow
  ?Frequency=A
  &Classification=HS
  &CommodityCodes=8542,8517,8471
  &Reporter=156
  &Partner=842
  &Flow=X
  &Period=2017,2018,2019,2020,2021,2022,2023,2024

导出文件：evidence/data/comtrade_chn_usa_electronics_2017-2024.csv

B. WITS Tariff
   evidence/data/wits_usa_tariff_on_chn_2017-2024.csv

C. USTR Section 301 列表
   List 1: 2018-07-06 生效, 500 亿美元
   List 2: 2018-08-23 生效, 160 亿美元
   List 3: 2018-09-24 生效, 2000 亿美元
   List 4A: 2019-09-01 生效

D. GTAP database 11 (2017 base)，聚合为 12 sectors x 4 regions

核验清单（7 项 — 必须线下完成）
- [ ] HS 8542、8517、8471 编码覆盖
- [ ] 单位一致（USD nominal vs 2017 实际）
- [ ] 镜像核验（China-as-Reporter vs USA-as-Reporter）
- [ ] 2024 年度发布滞后
- [ ] USTR 清单映射版本
- [ ] GTAP Electronics 部门口径
- [ ] 基年汇率
\end{lstlisting}

\subsection{第 4 步：/ueai-stage 02 — 模型与工具}

\noindent\textbf{用户输入}：
\begin{lstlisting}[style=prompt]
> /ueai-stage 02
> 我用 GTAP 11，做 25pp 关税冲击；本地路径 D:\GEMPACK12\GTAP\v11
\end{lstlisting}

\noindent\textbf{AI 实际写入 \texttt{02\_model/tariff\_shock\_25pp.cmf}}：
\begin{lstlisting}[style=md]
! 02_model/tariff_shock_25pp.cmf
! US +25pp tariff on imports of ELE from China, 2017 base, 12 x 4 agg
File (new) tariff25pp.sl4 ;
File (new) tariff25pp.har ;
File rwsbase = D:\GEMPACK12\GTAP\v11\base12x4.har ;
File rwspara = D:\GEMPACK12\GTAP\v11\para12x4.har ;
File gtapsets = D:\GEMPACK12\GTAP\v11\sets12x4.har ;
Verbal Description = "US +25pp tariff on ELE from CHN, 2017 base" ;
Auxiliary files = tariff25pp ;
Method = Gragg ;
Steps = 3 6 9 ;
Subintervals = 3 ;
Automatic accuracy = yes ;

Shock tms("ELE", "CHN", "USA") = 25 ;
\end{lstlisting}

\noindent\textbf{AI 同时写入 \texttt{02\_model/runGTAP.bat} 一键运行}：
\begin{lstlisting}[style=md]
:: 02_model/runGTAP.bat
@echo off
cd /d %~dp0
set GTAP=C:\GEMPACK12\GTAP\v11
runGTAP.exe tariff_shock_25pp.cmf > tariff25pp.log 2>&1
if %ERRORLEVEL% NEQ 0 (
    echo GTAP run failed. See tariff25pp.log
    pause
    exit /b %ERRORLEVEL%
)
runGTAP.exe tariff_shock_10pp.cmf > tariff10pp.log 2>&1
echo Done.
\end{lstlisting}

\noindent\textbf{学生在本机实际执行}：
\begin{lstlisting}[style=shell]
PS D:\courses\trade-db\2026-spring\student-007\tariff-impact\02_model>
  .\runGTAP.bat

Running GEMPACK ...
Begin solving model
Step 1 of 3 ... ok
Step 2 of 3 ... ok
Step 3 of 3 ... ok
Newton converged after 3 iterations
Wrote tariff25pp.sl4
Wrote tariff25pp.har
Done.
\end{lstlisting}

如果学生没有 GEMPACK，AI 同时给出 Stata DID 备选代码（位于 \texttt{02\_model/did\_section301.do}），见示范项目仓库。

\subsection{第 5 步：/ueai-stage 03 — 结果解释}

\noindent\textbf{用户输入}：
\begin{lstlisting}[style=prompt]
> /ueai-stage 03
> tariff25pp.sl4 + tariff10pp.sl4 已用 ViewSOL 导出到
> evidence/results/，关注中国 GDP、EV、电子部门出口
\end{lstlisting}

\noindent\textbf{AI 实际写入 \texttt{03\_results/result\_interpretation.md} 的关键数值段}：
\begin{lstlisting}[style=md]
关键数值（每条带 [来源: ...]）：

- 中国实际 GDP qgdp(CHN): -0.34 %
    [来源: evidence/results/tariff25pp_qgdp.csv]
- 中国福利 EV(CHN): -32.1 亿美元（2017 基价）
    [来源: evidence/results/tariff25pp_ev.csv]
- 中国电子部门对美出口 qxw("ELE",CHN,USA): -12.7 %
    [来源: evidence/results/tariff25pp_qxw_bilateral.csv]
- 中国电子部门总出口 qxw("ELE",CHN,ALL): -4.2 %
    [来源: evidence/results/tariff25pp_qxw_total.csv]
- 美国实际 GDP qgdp(USA): -0.11 %
- 美国福利 EV(USA): -18.6 亿美元（2017 基价）

异质性 - 受损前 3 部门（中国）

| 部门              | qo(CHN) | 主要原因          |
|-------------------|---------|-------------------|
| ELE（电子）       | -5.8 %  | 直接冲击部门      |
| MET（金属）       | -1.4 %  | 电子上游金属投入  |
| CHM（化工）       | -0.9 %  | 电子上游化工投入  |

稳健性：
1. 关税幅度 25pp → 10pp：中国 EV 缩为 -13.4 亿美元（线性缩减 ~60%）
2. 取消关税收入返还美国家庭：美国福利损失放大 18%
3. HS 章节扩展到 84+85+73+76：中国 EV 损失放大约 2.1 倍
4. DID 备选估计：-0.142 (log-points), cluster-SE = 0.038, p < 0.01

与文献对照：
- Caliendo & Parro (2015) 量级 0.08–0.45%，本研究 0.026% 在带内
- Amiti et al. (2019) 美国年化损失 ~168 亿，本研究 EV(USA) -18.6 亿
  （仅电子部门）量级吻合
- Itakura & Lee (2019) 全面贸易战中国 GDP 下降 0.5–0.8%，
  本研究单部门 -0.34% 在合理子集
\end{lstlisting}

\textbf{说明}：示范项目所引数值参考公开文献量级带，作为占位。真实学生项目须用本地 GTAP 求解的实际数值替换；每条 \texttt{[来源: ...]} 标记指向真实文件。

\subsection{第 6 步：/ueai-stage 04 — 论文与报告}

\noindent\textbf{用户输入}：
\begin{lstlisting}[style=prompt]
> /ueai-stage 04
\end{lstlisting}

\noindent\textbf{AI 实际写入 \texttt{04\_writing/paper\_outline.md} 的五维度大纲（节选）}：
\begin{lstlisting}[style=md]
## 1 引言
- 1.1 研究背景：中美贸易摩擦演进、Section 301 时间线
- 1.2 文献空白：现有 GTAP 多用 2014 GTAP 10；本文用 GTAP 11 (2017)
- 1.3 研究问题：2018 年以来美国对华加征关税对中国电子部门……
- 1.4 边际贡献（<= 3 条）
  - 用 GTAP 11 重新校准
  - CGE + DID 对照
  - 部门内异质性分解

## 2 文献综述
2.1 贸易战福利效应（Amiti et al. 2019; Fajgelbaum et al. 2020;
    Cavallo et al. 2021）
2.2 GTAP 应用（Caliendo–Parro 2015; Itakura–Lee 2019;
    Hertel et al. 2001）
2.3 中国视角（余淼杰 2019, 陈勇兵 2020, [ref待补]）

gap 句：现有研究对 (i) 2020 年后更新冲击、(ii) 电子子部门异质性、
(iii) 关税收入归宿假设的敏感性，缺乏分解。

## 3 数据与方法（引 01_data + 02_model）
## 4 结果（引 03_results，每数字带 [来源: ...]）
## 5 结论与政策含义（含"本研究不能说明…"段）

## 学术诚信声明
本研究的数据检索、模型设定、结果解释和论文写作过程中，使用了 AI
辅助工具 Claude Opus 4.7（通过 ueaiworkflow 工作流）。
AI 在 stage 00–04 提供了帮助；提示词、AI 反馈摘要与人工修改记录
见 05_feedback/ai_revision_log.md。
研究问题确立、数据核验、模型本地求解、结果数值判读和最终文本
由作者本人负责。
\end{lstlisting}

\subsection{第 7 步：/ueai-stage 05 — AI 反馈整理}

\begin{lstlisting}[style=prompt]
> /ueai-stage 05
\end{lstlisting}

\noindent\textbf{AI 整理 stage 00–04 全部 AI 段，写入 \texttt{05\_feedback/ai\_revision\_log.md}}：
\begin{lstlisting}[style=md]
| 阶段        | AI 模型          | 提示词来源                         | 状态                |
|-------------|------------------|------------------------------------|---------------------|
| 00_problem  | Claude Opus 4.7  | prompts/problem-framing.md         | 接受 + 学生补口径   |
| 01_data     | Claude Opus 4.7  | prompts/data-path-{comtrade,io}.md | 接受                |
| 02_model    | Claude Opus 4.7  | prompts/model-gtap-gempack.md      | 修改（本地路径）    |
| 03_results  | Claude Opus 4.7  | prompts/result-interpret.md        | 占位待替换          |
| 04_writing  | Claude Opus 4.7  | prompts/paper-structure.md         | 接受框架            |

使用合规自检
[✓] 所有 ## AI 辅助记录 段含时间 + 模型 + 提示词
[✓] 所有数字带 [来源: evidence/results/...]
[✓] 所有 ## 人工核验与修改 段保留待填位
[✓] 04_writing/paper_outline.md 含 ## 学术诚信声明 块
[✓] evidence/ 下文件由学生导出
\end{lstlisting}

\subsection{第 8 步：/ueai-review — 评阅模式}

\begin{lstlisting}[style=prompt]
> /ueai-review
> 评阅角色：paper-reviewer
\end{lstlisting}

\noindent\textbf{AI 追加到 \texttt{05\_feedback/ai\_revision\_log.md} 末尾的评阅表}：
\begin{lstlisting}[style=md]
## 评阅记录 — 2026-05-27 — paper-reviewer
模型：Claude Opus 4.7

| 维度        | 分数 | 反馈                                | 修改建议                       |
|-------------|------|-------------------------------------|--------------------------------|
| 1 研究问题  | 4/5  | Q 清晰，但边际贡献 3 在结果未体现  | §4.2 加 IC vs 通信 vs 计算机   |
| 2 文献逻辑  | 3/5  | 起点 6 条；gap 句完整              | 补 ≥ 4 条中文文献              |
| 3 数据方法  | 5/5  | 数据 / closure / 假设齐备          | 无                             |
| 4 结果解释  | 4/5  | 三件套齐；每数字带来源              | 数值仍占位；待本地替换         |
| 5 结论边界  | 5/5  | 含"本研究不能说明"段                | 无                             |

综合分：21 / 25（待数值替换后复评至 23–25）
\end{lstlisting}

\subsection{第 9 步：/ueai-archive — 归档清单}

\begin{lstlisting}[style=prompt]
> /ueai-archive
\end{lstlisting}

\noindent\textbf{AI 更新 \texttt{evidence/evidence\_checklist.md}}：
\begin{lstlisting}[style=md]
- [x] 1. 问题拆解：00_problem/problem_brief.md
- [~] 2. 数据路径：01_data + evidence/data/ (仅模板)
- [~] 3. 模型与工具：02_model + evidence/model/ (仅模板)
- [~] 4. 结果解释：03_results + evidence/results/ (占位)
- [x] 5. 论文与报告：04_writing/paper_outline.md
- [x] 6. AI 反馈：05_feedback/ai_revision_log.md
- [ ] 7. 最终报告或论文：evidence/paper/final.pdf

归档时间：2026-05-27
[Warning] 4/7 完整 + 3/7 占位 (示范预期)。
真实学生提交前须 >= 6/7 完整。
\end{lstlisting}

\subsection{第 10 步：/ueai-validate — 提交前完整性检查}

\begin{lstlisting}[style=prompt]
> /ueai-validate
\end{lstlisting}

\noindent\textbf{AI 输出报告（不写入文件，回显终端）}：
\begin{lstlisting}[style=shell]
ueai-validate report — tariff-impact — 2026-05-27

C1 [✓] Six stage files exist
C2 [✓] AI sections complete (5/5)
C3 [~] Human verification (0/5 — 示范项目预期)
C4 [✓] Evidence checklist exists
C5 [~] Evidence count (4/7 含占位)
C6 [✓] AI revision log non-empty
C7 [✓] No fabricated numbers (所有数字带 [来源: ...])
C8 [✓] Academic integrity statement present

VERDICT: NEEDS WORK (示范项目预期状态)

Required fixes (真实学生):
- 完成 5/5 的 "人工核验与修改" 段 (C3)
- 把 evidence/results/ 下的占位 .csv 替换为本地求解结果 (C5)
- 编译 LaTeX 生成 evidence/paper/final.pdf
- 复跑 /ueai-validate -> 预期 VERDICT: READY FOR SUBMISSION
\end{lstlisting}

%--------------------------------------------------------------------
\section{投入产出统计}
%--------------------------------------------------------------------

整套 \texttt{tariff-impact} 示范在项目负责人本机 Claude Code 中实际跑通；AI 端总计写入约 \textbf{11{,}000 字}；学生端实操工作量约 \textbf{7 小时}（GTAP 求解 + 数据导出占大头）。AI 没有替学生写一句正文——所有正文段落由学生在 \texttt{04\_writing/paper\_outline.md} 大纲基础上独立撰写。

\begin{center}
\small
\begin{tabularx}{\textwidth}{l c X r}
\toprule
\textbf{阶段命令} & \textbf{AI 写入} & \textbf{学生须完成} & \textbf{学生耗时} \\
\midrule
\texttt{/ueai-new}      & 8 个文件骨架    & 确认 slug 与课程        & 30 秒 \\
\texttt{/ueai-stage 00} & $\sim$800 字   & 5 项口径独立核验        & 30 分钟 \\
\texttt{/ueai-stage 01} & $\sim$1400 字  & UN Comtrade + WITS 导出 & 60 分钟 \\
\texttt{/ueai-stage 02} & $\sim$2200 字  & 本地跑 GTAP             & 90 分钟 \\
\texttt{/ueai-stage 03} & $\sim$2000 字  & 用本地数替换占位        & 60 分钟 \\
\texttt{/ueai-stage 04} & $\sim$2200 字  & 补文献 / 写摘要 / 写诚信声明 & 90 分钟 \\
\texttt{/ueai-stage 05} & $\sim$1100 字  & 校对 AI 日志            & 5 分钟 \\
\texttt{/ueai-review}   & $\sim$600 字   & 据修改建议改稿          & 90 分钟 \\
\texttt{/ueai-archive}  & $\sim$300 字   & 把缺的 evidence 补齐    & 15 分钟 \\
\texttt{/ueai-validate} & $\sim$250 字   & 修补失败项              & 10 分钟 \\
\midrule
\textbf{合计}            & $\sim$11{,}000 & ---                     & $\sim$7 小时 \\
\bottomrule
\end{tabularx}
\end{center}

%--------------------------------------------------------------------
\section{评价、归档与推广}
%--------------------------------------------------------------------

\subsection{推广应用价值}

工作流的核心是 ``\emph{把研究过程拆为六阶段、每阶段配 AI 技能、每阶段留人工核验段、每阶段沉淀证据}'' 这一通用结构。该结构可推广至：

\begin{itemize}[leftmargin=*]
  \item \textbf{计量经济学} —— 把 \texttt{prompts/model-regression.md} 升级为 DID / 工具变量 / RDD 系列；
  \item \textbf{国际贸易实务} —— 把 \texttt{prompts/data-path-comtrade.md} 扩展到 ITC Trademap / CEPII 等；
  \item \textbf{产业经济学 / 数字经济专题} —— 替换 \texttt{prompts/} 中的数据库与模型部分；
  \item \textbf{经贸研究与论文写作} —— 直接复用全部工作流，仅替换课程名与示范项目。
\end{itemize}

教师只需要在课程开始时把仓库 URL 发给学生，全班的目录结构、AI 使用披露、评阅记录就一次性统一了；学期末把每位学生 \texttt{<slug>/} 文件夹打包，即得到课程档案与项目结项佐证。

\subsection{对学生的价值}

学生在完成一次完整工作流后，自然形成对``经济研究是怎么一步步做出来的''的整体认识：从一个模糊兴趣到一组核验过的数据，再到一份带证据的论文。AI 在每一步给的帮助都被记录、被分类、被复核，因此 AI 既不会被滥用为替写工具，也不会被神秘化为黑箱。

\subsection{对教师的价值}

教师从评阅一份不可追溯的 PDF 转为评阅一组结构化的过程材料：数据来源截图、模型运行日志、AI 反馈日志、人工修改痕迹、五维度大纲、归档清单。期末成绩与项目结项佐证因此可以建立在过程证据上，而不是仅凭文本。

\subsection{引用}

\begin{lstlisting}[style=md]
ueaiworkflow [CP/OL]. 2026. https://github.com/ALLBLUEUK/ueaiworkflow
\end{lstlisting}

\subsection{许可}

MIT，见仓库根目录 \texttt{LICENSE}。

\end{document}
